\documentstyle[% 


righttag% 


,11pt% 


,amssymb% 


,russian%               remove in Eng. 


,izvru%                 Set izven% in Eng. 


]{amsart} 


 


% 


 


\newcommand{\sgn}{\operatorname{sgn}} 


\newcommand{\ctg}{\operatorname{ctg}} 


\newcommand{\arcctg}{\operatorname{arcctg}} 


\newcommand{\tts}[1]{} 


 


\theoremstyle{definition} 


\theorembodyfont{\rm} 


\newtheorem{cond}{\hskip\parindent Условие} 


 


%\hoffset=-15mm%                  For Eng. 


\setcounter{page}{3} 


\god{2004} 


\nomer{4~(503)} %                        Remove in Eng. 


%\nomer{Vol.\,48, No.\,4}%               Set in Eng. 


%\str{\pageref{begin}--\pageref{end}}%  Set in Eng. 


\udk{517.958} 


 


\author{\it Г.Н.\,Аглямзянова, Р.С.\,Хайруллин} 


% NB:   ^^ remove in Eng. 


\title{Задача Трикоми в классе функций,\\ 


неограниченных на характеристике} 


 


\date{} 


%\pagestyle{myheadings}%         Set in Eng. 


\pagestyle{plain} %              Remove in Eng. 


\begin{document} 


%\thispagestyle{plain}%          Set in Eng. 


\maketitle 


\label{begin} 


 


 


 


Рассмотрим уравнение 


\begin{equation} 


\sgn yu_{xx}+u_{yy}+\frac{2q}{y}u_x+\frac{2p}{y}u_y=0,\quad 


p<0, 


\end{equation} 


где $p$, $q$ --- вещественные параметры, в смешанной области $D$, 


эллиптическая 


подобласть которой $D_1$ совпадает со всей верхней полуплоскостью, а 


гиперболическая подобласть $D_2$ представляет собой треугольник, 


ограниченный характеристиками $BC$: $x-y=1$, $AB$: $x+y=0$ и отрезком $AC$ 


оси абсцисс. 


 


Введем следующие обозначения: $2p=\beta+\beta'$, $2q=\beta-\beta'$. Задача 


Трикоми для уравнения (1) 


была исследована в работе [1]. Было доказано, что в случае 


$\beta>0$, $\beta'<0$ задача имеет единственное решение при выполнении 


$n+1$ или $n+2$ 


условий разрешимости в зависимости от значения параметра $c$, входящего в 


условие склеивания. Здесь $n\in N_0=N\cup\{0\}$ и выполняется неравенство 


$0<\beta-n<1$, $\beta_0=\beta-n$. 


 


Отличительной особенностью данной статьи является исследование 


задачи Трикоми в классе функций, неограниченных на характеристике $BC$, что 


позволяет снять указанные условия разрешимости. 


 


Пусть $k,m\in N$ и $0<\beta'+k<1$, $\beta'_0=\beta'+k$, $0<2p+m<1$; 


$\delta=2p+m$. 


Будем исследовать случай $1<\beta_0+\beta'_0<2$. Тогда имеем 


$\beta_0>\delta$, $\beta'_0>\delta$, $m=k-n-1$. 


\medskip 


 


{\bf Задача T.} {\it Найти функцию $u(x,y)$ со свойствами 


\begin{itemize} 


\item[1)] $u(x,y)$ принадлежит $C(D_1\cup D_2\cup AB\cup\{(x;0)\})$ и 


ограничена на бесконечности\/$;$ 


\item[2)] $u(x,y)$ принадлежит $C^2(D_1\cup D_2)$ и удовлетворяет уравнению 


$(1)$ в $D_1\cup D_2;$ 


\item[3)] существуют пределы 


\begin{equation} 


v_j(x)=\lim\begin{Sb}y\to0\\ (x,y)\in D_j\end{Sb}|y|^{2p} 


[u(x,y)-A^j_{p,q}(x,y,\tau)]_y,\quad 0<x<1, \quad j=1,2,


\end{equation} 


и на $AC$ выполняется условие склеивания 


\begin{equation} 


v_1(x)=cv_2(x),\quad 0<x<1, 


\end{equation} 


где 


\begin{equation} 


\tau(x)=u(x,0),\quad 0\le x\le 1, 


\end{equation} 


$A^j_{p,q}(x,y,\tau)=\sum\limits^m_{l=1} \frac{a^j_l}{l!} 


\tau^{(l)}(x)y^l$, $a^1_l=l_1(-1)^l\int\limits^\pi_0 e^{2q\xi} \cos^l 


\xi\sin^{-2p-l}\xi\,d\xi$, $l_1=B(1-p-iq,\;1-p+iq)/4^p \pi e^{\pi q}$, $B$ 


--- бета-функция, $c>0$ --- параметр, 


$a^2_l=F(-l,\beta,\beta+\beta',2)$, $F$ --- гипергеометрическая функция 


Гаусса\/$;$ 


\item[4)] $u(x,y)$ удовлетворяет краевым условиям 


\begin{center}


$u(x,0)=0$, \ $x\le0$ \ или \ $x\ge1$, \ $u(x,y)|_{AB}=\psi(x)$, \ $0\le 


x\le\frac{1}{2}$. 


\end{center}


\end{itemize} 


} 


 


\noindent{}Здесь $\psi(x)$ --- заданная функция, для которой выполняется 


 


\begin{cond} 


Функция $\psi\big(\frac{x}{2}\big)\in C[0,1)$ при $x=0$ имеет нуль порядка 


больше $1-\delta$, 


а при $x=1$ имеет особенность порядка меньше $\beta$. 


\end{cond} 


 


 


Решение задачи Трикоми будем искать в классе функций $\tau(x)$, $v_i(x)$, 


для которых выполняется 


 


\begin{cond} 


Функция $\tau(x)$ принадлежит $C[0,1]\cap C^{m-1}(0,1]\cap C^k(0,1)$,


$\tau^{(m)}(x)$ при $x=1$ имеет 


особенность порядка меньше $\delta$, $\tau(x)$ при $x=0$ имеет нуль порядка 


больше 


$1-\delta$; функции $v_i(x)\in C^n(0,1)$ могут иметь особенности при $x=1$ 


порядка 


меньше~1, при $x=0$ --- порядка меньше $m$. 


\end{cond} 


 


Задачу будем решать методом интегральных уравнений. Соотношения из 


эллиптической и гиперболической подобластей, полученные в [1], 


справедливы для класса функций, где $\tau(x)\in C[0,1]\cap C^m(0,1]$, 


$v_i(x)\in C^n(0,1)$. Поэтому эти 


соотношения можно использовать и в данном случае. Они имеют вид 


\begin{align} 


v_1(x)&=-l_1\Gamma(2p)D^{1-2p}_{0x}\tau(x)-l_1e^{2\pi q} 


\Gamma(2p)D^{1-2p}_{x1}\tau(x),\\ 


% 


v_2(x)&=l_2\Gamma(2p)D^{1-2p}_{0x}\tau(x)- 


\frac{\Gamma(1-\beta')2^{1-\beta-\beta'}}{\Gamma(1-\beta-\beta')} 


x^{\beta'}D^{1-\beta}_{0x}\psi\bigg(\frac{x}{2} \bigg), 


\end{align} 


где $l_2=\frac{2\Gamma(1-\beta')}{\Gamma(\beta)(1-\beta-\beta')}$, 


$D^{1-2p}_{0x}$, $D^{1-2p}_{x1}$ --- дробные производные. 


 


Из условий 1, 2 и формулы (5) получим


\begin{equation} 


\tau^{(s)}(1)=0,\quad s=\ov{0,m-1}. 


\end{equation} 


 


Перейдем к выводу интегрального уравнения. Используя условие склеивания 


(3), из соотношений (5) и (6) имеем 


\begin{equation} 


(l_1+cl_2)D^{1-2p}_{0x}\tau(x)+l_1e^{2\pi q}D^{1-2p}_{x1}\tau(x) 


=l_3x^{\beta'}D^{1-\beta}_{0x}\psi\bigg(\frac{x}{2} \bigg), 


\end{equation} 


где $l_3=\frac{c2^{1-\beta-\beta'}\Gamma(1-\beta')} 


{\Gamma(2p)(1-\beta-\beta')}$. 


 


Дальнейший способ вывода уравнения отличается от предложенного в [1]. Обе 


части равенства (8) при $x=0$ имеют особенности порядка меньше $m$. 


Применим к нему оператор $\int\limits^x_0(x-\xi)^{-\delta} \xi^m\dotsc 


d\xi$. Используя формулу ([2], с.\,78) 


$x^{-a-b}M(a,b,l,x)=\frac{d^l}{dx^l}x^{l-a-b}M(a,b-l,0,x)$, где 


$M(a,b,l,x)=\int\limits^x_0 g^{(l)}(\sigma)\sigma^b 


(x-\sigma)^{a-1}d\sigma$, $a>0$, $b>l-1$, 


получим 


\begin{multline} 


(l_1+cl_2)x^{1-2p}\frac{d^m}{dx^m}x^{\delta-1} 


D^{\delta-1}_{0x}D^{1-\delta}_{0x}\tau(x)+l_1e^{2\pi q} 


x^{1-2p}\frac{d^m}{dx^m}x^{\delta-1}D^{\delta-1}_{0x} 


D^{1-\delta}_{x1}\tau(x)=\\ 


% 


=l_3\int^x_0(x-\xi)^{-\delta}\xi^{m+\beta'}D^{1-\beta}_{0\xi} 


\psi\bigg(\frac{\xi}{2} \bigg)d\xi. 


\end{multline} 


На основе формул композиции дробных производных и интегралов ([3], 


сс.\,18, 24) выражение (9) примет вид 


\begin{multline*} 


(l_1+cl_2)x^{1-2p}\frac{d^m}{dx^m}x^{\delta-1}\tau(x) 


+l_1e^{2\pi q}x^{1-2p}\frac{d^m}{dx^m} 


x^{\delta-1}\bigg[\tau(x)\cos\pi(1-\delta)-\\ 


% 


-\frac{\sin\pi(1-\delta)}{\pi}\int^1_0\bigg(\frac{x}{\sigma} \bigg)^{1-\delta} 


\frac{\tau(\sigma)d\sigma}{\sigma-x}\bigg]=l_3D^{\delta-1}_{0x} 


x^{m+\beta'}D^{1-\beta}_{0x}\bigg(\frac{x}{2} \bigg). 


\end{multline*} 


Обозначим $\mu(x)=x^{\delta-1}\tau(x)$ и после приведения подобных членов 


получим интегральное уравнение 


\begin{equation} 


\mu^{(m)}(x)\ctg\pi\varphi-\frac{1}{\pi}\frac{d^m}{dx^m} 


\int^1_0\frac{\mu(\sigma)d\sigma}{\sigma-x}=f(x), 


\end{equation} 


где 


$$ 


f(x)=\frac{l_3\Gamma(1-\delta)x^{2p-1}} 


{l_1e^{2\pi q}\sin\pi(1-\delta)}D^{\delta-1}_{0x} 


x^{m+\beta'}D^{1-\beta}_{0x}\psi\bigg(\frac{x}{2} \bigg),\quad 


\varphi=\frac{1}{\pi}\arcctg\bigg(\frac{l_1+cl_2} 


{l_1e^{2\pi q}\sin\pi(1-\delta)}+\ctg\pi(1-\delta)\bigg). 


$$ 


 


У функции $\mu^{(m)}(x)$ допускаются особенности при $x=0$ порядка меньше 


$m$, 


при $x=1$ --- меньше $\delta$. Такие же особенности имеет функция $f(x)$. 


 


Учитывая равенства (7), получим 


\begin{equation} 


\mu^{(s)}(1)=0,\quad s=\ov{0,m-1}. 


\end{equation} 


Преобразуем уравнение (10). Будем заносить операцию дифференцирования 


под знак интеграла. Первую производную можно вычислить, непосредственно 


дифференцируя плотность интеграла типа Коши. Для следующих производных 


можно использовать формулу дифференцирования интеграла с ядром типа Коши 


с плотностью, допускающей интегрируемые особенности. Она имеет вид 


[4] 


$$ 


Q(x)\frac{d}{dx}\int^1_0\frac{g(\sigma)d\sigma}{\sigma-x}= 


\int^1_0\frac{[Q(\sigma)g(\sigma)]'}{\sigma-x}d\sigma+ 


\int^1_0\frac{g(\sigma)}{\sigma-x} 


\frac{Q(x)-Q(\sigma)}{\sigma-x}d\sigma. 


$$ 


 


Чтобы избавиться от особенности при $x=0$, положим $Q=x$. Тогда 


последнее соотношение примет вид 


\begin{equation} 


\frac{d}{dx}\int^1_0\frac{g(\sigma)d\sigma}{\sigma-x}= 


\int^1_0\bigg(\frac{\sigma}{x} \bigg)\frac{g'(\sigma)}{\sigma-x}d\sigma. 


\end{equation} 


Рассмотрим выражение 


\begin{multline*} 


\frac{d^{m-s}}{dx^{m-s}}\int^1_0\bigg(\frac{\sigma}{x} \bigg)^{s-1} 


\frac{\mu^{(s)}(\sigma)d\sigma}{\sigma-x}=\\ 


% 


=\frac{d^{m-s-1}}{dx^{m-s-1}}\bigg\{-\frac{s-1}{x^s} 


\int^1_0\frac{\sigma^{s-1}\mu^{(s)}(\sigma)}{\sigma-x}d\sigma+ 


\frac{1}{x^{s-1}}\frac{d}{dx}\int^1_0 


\frac{\sigma^{s-1}\mu^{(s)}(\sigma)}{\sigma-x}d\sigma\bigg\},


\quad s=\ov{0,m-1}. 


\end{multline*} 


Применяя формулу (12) ко второму слагаемому в правой части, имеем 


\begin{gather*} 


\frac{d^{m-s}}{dx^{m-s}}\int^1_0\bigg(\frac{\sigma}{x} \bigg)^{s-1} 


\frac{\mu^{(s)}(\sigma)d\sigma}{\sigma-x}=\frac{d^{m-s-1}} {dx^{m-s-1}} 


\bigg\{\frac{-(s-1)}{x^s}\int^1_0 


\frac{\sigma^{s-1}\mu^{(s)}(\sigma)}{\sigma-x}d\sigma+ \\ 


% 


+\frac{1}{x^{s-1}}\int^1_0\frac{\sigma}{x} 


\frac{(s-1)\sigma^{s-2}\mu^{(s)}(\sigma)}{\sigma-x}d\sigma+ 


\frac{1}{x^{s-1}}\int^1_0\frac{\sigma}{x} 


\frac{\sigma^{s-1}\mu^{(s+1)}(\sigma)}{\sigma-x}d\sigma\bigg\}= 


\frac{d^{m-s-1}}{dx^{m-s-1}}\int^1_0\bigg(\frac{\sigma}{x} \bigg)^s 


\frac{\mu^{(s+1)}(\sigma)d\sigma}{\sigma-x}. 


\end{gather*} 


Отсюда получим 


$$ 


\frac{d^{m-1}}{dx^{m-1}}\int^1_0\frac{\mu(\sigma)d\sigma}{\sigma-x} 


=\int^1_0\bigg(\frac{\sigma}{x} \bigg)^{m-1} 


\frac{\mu^{(m)}(\sigma)d\sigma}{\sigma-x}. 


$$ 


В результате уравнение (10) примет вид 


\begin{equation} 


\mu^{(m)}(x)\ctg\pi\varphi-\frac{1}{\pi}\int^1_0 


\bigg(\frac{\sigma}{x} \bigg)^{m-1} 


\frac{\mu^{(m)}(\sigma)d\sigma}{\sigma-x}=f(x). 


\end{equation} 


 


Обозначим $\rho(x)=\mu^{(m)}(x)x^{m-1}$, $g(x)=x^{m-1}f(x)$. Тогда (13) 


перейдет в уравнение 


\begin{equation} 


\rho(x)\ctg\pi\varphi-\frac{1}{\pi}\int^1_0 


\frac{\rho(\sigma)d\sigma}{\sigma-x}=g(x). 


\end{equation} 


 


Функция $g(x)$ может иметь особенности при $x=0$ порядка меньше 1, при 


$x=1$ --- меньше~$\delta$. Такие же особенности имеет функция $\psi(x)$. 


Заметим, что 


уравнение (14) в классе функций, неограниченных при $x=0$ и $x=1$, 


имеет решение с особенностью при $x=0$ порядка $\varphi$, при $x=1$ --- 


порядка $1-\varphi$ ([5]). 


Но $\varphi<\frac{1}{\pi}\arcctg(\ctg\pi(1-\delta))=1-\delta$, 


т.\,е. $1-\varphi>\delta$. Поэтому необходимо использовать 


формулу решения, ограниченного при $x=1$. Такое решение единственно. Далее 


решение исходной задачи находится по стандартной схеме. При 


восстановлении функции $\mu(x)$ используются равенства (11). Единственность 


решения задачи Трикоми следует из однозначности определения 


вспомогательных функций 


и единственности решений задачи Дирихле и задачи типа Коши. 


 


Итак, доказана 


 


\begin{thm} 


Если функция $\psi(x)$ удовлетворяет условию $1$, то в классе функций 


$\tau(x)$, $v_i(x)$, удовлетворяющих условию $2$, задача Трикоми имеет 


единственное решение.


\end{thm} 


 


Исследуем поведение решения $u(x,y)$ на характеристике $BC$. Решение задачи 


типа Коши (1), (2), (4) записывается формулой [1] 


\begin{multline*} 


u(x,y)=\sum^{k-1}_{s=0}\frac{\tau^{(s)}(x-y)2^s(\beta')_s} 


{s!(\beta+\beta')_s}y^s-\sum^{n-1}_{s=0} 


\frac{v^{(s)}(x-y)2^s(-1)^s(1-\beta)_s} 


{s!(1-\beta-\beta')_{s+1}}(-y)^{1-\beta-\beta'+s}+ \\ 


% 


+\frac{\Gamma(\beta+\beta')2^ky^k} 


{\Gamma(\beta')\Gamma(\beta+k)}\int^1_0\tau^{(k)}(\varsigma) 


\sigma^{\beta+k-1}F(\beta,\,1-\beta',\,\beta+k,\,\sigma)d\sigma-\\ 


% 


-\frac{\Gamma(1-\beta-\beta')2^n(-1)^n(-y)^{1-\beta-\beta'+n}} 


{\Gamma(1-\beta'+n)\Gamma(1-\beta)}\int^1_0 v^{(n)} 


(\varsigma)\sigma^{n-\beta'}F(\beta,\,1-\beta',\,1-\beta'+n,\,\sigma) 


d\sigma. 


\end{multline*} 


С новой переменной интегрирования $\varsigma=x+y(1-2\sigma)$ получим 


\begin{multline} 


u(x,y)=\sum^{k-1}_{s=0}\frac{\tau^{(s)}(x-y)2^s(\beta')_s} 


{s!(\beta+\beta')_s}y^s-\sum^{n-1}_{s=0} 


\frac{v^{(s)}(x-y)2^s(-1)^s(1-\beta)_s} 


{s!(1-\beta-\beta')_{s+1}}(-y)^{1-\beta-\beta'+s}+\\ 


% 


+\frac{\Gamma(\beta+\beta')2^{-\beta}y^{-\beta}} 


{\Gamma(\beta')\Gamma(\beta+k)}\int^{x-y}_{x+y}\tau^{(k)}(\varsigma) 


(x+y-\varsigma)^{\beta+k-1}F\bigg(\beta,\,1-\beta',\,\beta+k,\, 


\frac{x+y-\varsigma}{2y}\bigg)d\varsigma-\\ 


% 


-\frac{\Gamma(1-\beta-\beta')2^{\beta'-1}(-1)^{\beta'-1}(-y)^{-\beta}} 


{\Gamma(1-\beta'+n)\Gamma(1-\beta)}\int^{x-y}_{x+y} 


v^{(n)}(\varsigma)(x+y-\varsigma)^{n-\beta'}F\bigg(\beta,\, 


1-\beta',\,1-\beta'+n,\,\frac{x+y-\varsigma}{2y}\bigg)d\varsigma=\\ 


% 


=S_1(x-y)-S_2(x-y)+I_1(x-y)-I_2(x-y). 


\end{multline} 


Оценим каждое слагаемое из правой части (15) при $x-y=t\to1$. Учитывая 


\begin{equation} 


\tau^{(k-1)}(t)=o\big((1-t)^{1-2p-k}\big)\ \ 


\text{при} \ \ t\to1, 


\end{equation} 


имеем 


$$ 


|S_1(t)|\le\sum^{k-1}_{s=0}\bigg| 


\frac{\tau^{(s)}(x-y)2^s(\beta')_s} 


{s!(\beta+\beta')_s}(x-t)^s\bigg|\le 


\frac{2^{k-1}(\beta')_{k-1}}{(k-1)!(\beta+\beta')_{k-1}} 


M(x-t)^{k-1}(1-t)^{-\delta-n}. 


$$ 


Отсюда получим $S_1=o\big((1-t)^{-n}\big)$ при $t\to1$. 


 


Аналогично из $v(t)=o\big((1-t)^{-1}\big)$ при $t\to1$ имеем 


$S_2(t)=o\big((1-t)^{-n}\big)$ при $t\to1$. 


 


С использованием переменной $t=x-y$ запишем 


$$ 


I_1(t)=\frac{\Gamma(\beta+\beta')2^{-\beta}(t-x)^{-\beta}} 


{\Gamma(\beta')\Gamma(\beta+k)}\int^t_{2x-t}\tau^{(k)} 


(\varsigma)(\varsigma-(2x-t))^{\beta+k-1}F\bigg(\beta,\,1-\beta', 


\,\beta+k,\,\frac{\varsigma-(2x-t)}{2(t-x)}\bigg)d\varsigma. 


$$ 


 


Применяя последовательно соотношение (16) и формулу автотрансформации 


для гипергеометрической функции, получим 


\begin{multline*} 


|I_1(t)|\le C_1(t-x)^{-\beta}\int^t_{2x-t}(1-\zeta)^{-2p-k+\varepsilon} 


(\varsigma-(2x-t))^{\beta+k-1}\bigg(1-\frac{\varsigma-(2x-t)}{2(t-x)} 


\bigg)^{\beta'_0-1}d\varsigma=\\ 


% 


=C_1(t-x)^{1-\beta-\beta'_0}\int^t_{2x-t}(1-\zeta)^{-2p-k+\varepsilon} 


(\varsigma-(2x-t))^{\beta+k-1}(t-\zeta)^{\beta'_0-1}d\varsigma, 


\end{multline*} 


где $\zeta=t-(2t-2x)z$. Тогда 


$$ 


|I_1(t)|\le C_3(t-x)^k(1-t)^{-2p-k+\varepsilon} 


F(\beta'_0,\,2p+k-\varepsilon,\,\beta'_0+\beta+k,\,\eta), 


\quad \eta=2(x-t)/(1-t). 


$$ 


Используя формулу Больца (15.3.4) [6], имеем 


\begin{gather*} 


|I_1(t)|\le C_3(t-x)^k(1-t)^{-2p-k+\varepsilon} 


(1-\eta)^{-\beta'_0}F\bigg(\beta'_0,\,\beta'_0+\beta-2p+\varepsilon, 


\,\beta'_0+\beta+k,\,\frac{\eta}{\eta-1}\bigg)=\\ 


% 


=C_3(t-x)^k(1-t)^{-2p-k+\beta'_0+\varepsilon} 


(1-(2x-t))^{-\beta'_0}F\bigg(\beta'_0,\, 


\beta'_0+\beta-2p+\varepsilon,\,\beta'_0+\beta+k,\, 


\frac{2x-2t}{2x-t-1}\bigg)=\\ 


=O\big((1-t)^{-\beta+\varepsilon}\big)\ \ 


\text{при}\ \ t\to1. 


\end{gather*} 


 


Аналогично получим $I_2(t)=O\big((1-t)^{-\beta+\varepsilon}\big)$ при 


$t\to1$. 


 


Таким образом, доказана 


 


\begin{thm} 


Решение $u(x,y)$ при $x-y=t\to1$ имеет особенность порядка ниже $\beta$. 


\end{thm} 


 


 


\begin{thebibliography}{99} 


 


\bibitem{1} 


Хайруллин Р.С. {\em Задача Трикоми для одного уравнения с сингулярными 


коэффициентами} // Изв. вузов. Математика. -- 1996. -- \No\,3. 


-- С.\,68--76. 


\bibitem{2} 


Хайруллин Р.С. {\em Задача Трикоми для дифференциальных уравнений с 


сильным вырождением на линии изменения типа}: Дисс. $\dotsc$ докт. 


физ.-матем. наук. -- Казань, 1994. -- 289~с. 


\bibitem{3} 


Смирнов М.М. {\em Уравнения смешанного типа}. -- М.: Наука, 1970. -- 295~с. 


\bibitem{4} 


Чибрикова Л.И., Плещинский Н.Б. {\em Об интегральных уравнениях с 


обобщенными логарифмическими и степенными ядрами} // Изв. вузов. 


Математика. -- 1978. -- \No\,6. -- С.\,129--146. 


\bibitem{5} 


Гахов Ф.Д. {\em Краевые задачи}. -- 3-е изд. -- М.: Наука, 1977. -- 640~с. 


\bibitem{6} 


{\em Справочник по специальным функциям}: Пер. с англ. / Под ред. 


М.\,Абрамовица, И.\,Стиган. -- М.: Наука, 1979. -- 830~с. 


 


\end{thebibliography} 


 


\vskip 2cm 


 


\post{\it Казанская государственная}{\it Поступила} 


\post{\it архитектурно-строительная академия}{09.07.2002} 


 


\label{end} 


\end{document} 


